\documentclass[11pt]{article}

%%METADATA
\title{GA1025 Macro Theory \\Solution to Problem Set 1}
\author{
Junbiao Chen\thanks{E-mail: jc14076@nyu.edu.}
}
\date{\today}


%%PACKAGES
\usepackage{mdframed} % For boxed environments
\usepackage{graphicx}
\usepackage{grffile}
\usepackage{tabularx}
\usepackage{setspace}
\usepackage{amsmath,amsthm,amssymb}
\usepackage[hyphens]{url}
\usepackage{natbib}
\usepackage[font=normalsize,labelfont=bf]{caption}
\usepackage[margin=1in]{geometry}
\usepackage{hyperref}
\hypersetup{colorlinks=true,urlcolor=blue,citecolor=blue}
\usepackage{stmaryrd}  %Package with \boxast command
\usepackage{enumerate}% http://ctan.org/pkg/enumerate %Supports lowercase Roman-letter enumeration
\usepackage{verbatim} %Package with \begin{comment} environment
%\usepackage{enumitem}
\usepackage{physics}
\usepackage{tikz}
\usepackage{listings}
\usepackage{upquote}
\usepackage{booktabs} %Package with \toprule and \bottomrule
\usepackage{etoc}     %Package with \localtableofcontents
\usepackage{placeins}    %Package that prevent repositioning the tables
\usepackage{multicol}
\usepackage{bm}
\usepackage{subfig}
\usepackage{csquotes}

\definecolor{dkgreen}{rgb}{0,0.6,0}
\definecolor{gray}{rgb}{0.5,0.5,0.5}
\definecolor{mauve}{rgb}{0.58,0,0.82}

\lstset{language=bash,
  frame=tb,
  aboveskip=3mm,
  belowskip=3mm,
  showstringspaces=false,
  columns=flexible,
  basicstyle={\small\ttfamily},
  numbers=none,
  numberstyle=\tiny\color{gray},
  keywordstyle=\color{blue},
  commentstyle=\color{dkgreen},
  stringstyle=\color{mauve},
  breaklines=true,
  breakatwhitespace=false,
  tabsize=3
}

\lstset{language=C,
  aboveskip=3mm,
  belowskip=3mm,
  showstringspaces=false,
  columns=flexible,
  basicstyle={\small\ttfamily},
  numbers=none,
  numberstyle=\tiny\color{gray},
  keywordstyle=\color{blue},
  commentstyle=\color{dkgreen},
  stringstyle=\color{mauve},
  breaklines=true,
  breakatwhitespace=false,
  tabsize=4
}

\definecolor{lightblue}{rgb}{0.68, 0.85, 0.9} %

%CUSTOM DEFINITIONS
\theoremstyle{definition}
\newtheorem{definition}{Definition}[section]
\newtheorem*{remark}{Remark}
\setcounter{secnumdepth}{3}
\usepackage{tikz}
\usetikzlibrary{arrows.meta}
\usetikzlibrary{automata, positioning, arrows, calc}

\tikzset{
	->,  % makes the edges directed
	>=stealth, % makes the arrow heads bold
	shorten >=2pt, shorten <=2pt, % shorten the arrow
	node distance=3cm, % specifies the minimum distance between two nodes. Change if n
	every state/.style={draw=blue!55,very thick,fill=blue!20}, % sets the properties for each ’state’ n
	initial text=$ $, % sets the text that appears on the start arrow
 }

%% PROPOSITION
% Define the Proposition environment
\newmdenv[
  innerleftmargin=10pt, 
  innerrightmargin=10pt,
  innertopmargin=10pt,
  innerbottommargin=10pt,
  linecolor=black, 
  linewidth=1pt,
  backgroundcolor=white, 
  roundcorner=5pt
]{propositionbox}

\newtheoremstyle{boldtitle} % Define a new theorem style
  {10pt} % Space above
  {10pt} % Space below
  {\itshape} % Body font
  {} % Indent amount
  {\bfseries} % Theorem head font
  {.} % Punctuation after theorem head
  { } % Space after theorem head
  {} % Theorem head spec

\theoremstyle{boldtitle} % Use the custom style
\newtheorem{proposition}{Proposition} % Define the proposition environment

% Redefine the proposition environment to use the box
\newenvironment{boxedproposition}[1][]
{\begin{propositionbox}\begin{proposition}[#1]}
{\end{proposition}\end{propositionbox}}

%%FORMATTING
\usepackage[bottom]{footmisc}
\onehalfspacing
\numberwithin{equation}{section}
\numberwithin{figure}{section}
\numberwithin{table}{section}
\bibliographystyle{../bib/aeanobold-oxford}


%main text
\begin{document}
\section{Search Models}
\subsection{Exercise 7.1}
The Bellman equation for a McCall-model with a chance of unemployment is as follows
\begin{equation}
    \begin{aligned}
    \bar{V} & = c + \beta \bigg( \phi \bar{V} + (1 - \phi) \mathbb{E} \left[ v(w')\right]  \bigg) \\
    v(w) & = \max_{a \in \{ \text{Accept, Reject} \}} \{ \frac{w}{1 - \beta}, \beta (\phi (c + \bar{V}) + (1 - \phi) \mathbb{E} \left[ v(w') \right]  ) \}
    \end{aligned}
\end{equation}
where $\bar{V}$ denotes the continuation value of being unemployed.
$v(w)$ denotes the value function conditional on receiving an offer.
$\mathbb{E}\left[ v(w') \right] = \int_{0}^B v(w') dF(w')$ denotes the expected value conditional on receiving an offer.


\subsection{Exercise 7.2}
\textbf{A McCall Model with two offers per period.}
Since $F(x) = \text{Pr}(z \leq x)$,  
the cumulative distribution function of the maximum of two random variables i.i.d. drawn from $F(x)$ is  
$F_{\max}(x) = \text{Pr}(x_1 < x) \cdot \text{Pr}(x_2 < x) = [F(x)]^2$.  

\begin{equation}
    v(w) = \max_{a \in \{\text{Accept, Reject\}}}  
    \bigg( \frac{w}{1 - \beta}, c + \beta \mathbb{E}_{\max}[v(w')] \bigg)
\end{equation}

where $\mathbb{E}_{\max}[v(w')] = \int_0^{B} v(w') \, d[F(w')]^2$.  
Since $\mathbb{E}_{\max}[v(w')] = \int_0^{B} v(w') \, d[F(w')]^2  
= 2 \int_0^{B} v(w') f(w') \, dw' = 2 \mathbb{E}[v(w')]$,  
we have  
\[
\bar{w}_{\max} = (1 - \beta) \big(c + \beta \mathbb{E}_{\max}[v(w')]\big)  
> (1 - \beta) \big(c + \beta \mathbb{E}[v(w')]\big) = \bar{w}.
\]  
\(\blacksquare\)

\subsection{Exercise 7.3}
\textbf{A McCall Model with $n$ iid offers per period.}
Generalize the formulation in Exercise 7.2, we have
\begin{equation}
    v(w) = \max_{a \in \{\text{Accept, Reject} \} }  \bigg( \frac{w}{1 - \beta}, c + \beta \mathbb{E}_\max [v(w')] \bigg)
\end{equation}
where $\mathbb{E}[v(w')] = \int_0^{B} \underbrace{\int_0^{w'} \dots \int_0^{w'}}_{n-1} v(w') \underbrace{dF(w') \dots dF(w')}_{n}$
$\mathbb{E}_\max [v(w')] = n \int_{0}^B v(w') f(w') dw'$.

%% =========== %% ========= %%
\bibliography{../bib/notes.bib}

\end{document}